\documentclass{article}

\usepackage[utf8]{inputenc}
\usepackage{helvet}
\usepackage{geometry}
\usepackage{xcolor}
\usepackage{eso-pic}
\usepackage{fancyhdr}
\usepackage{parskip}
\usepackage{microtype}
\usepackage[english, ukrainian]{babel}
\usepackage{url}
\usepackage{amsmath}
\usepackage{amssymb}
\usepackage{amsthm}
\usepackage{longtable}
\usepackage{booktabs}
\usepackage{hyperref}
\usepackage{titlesec}

\definecolor{nato}{HTML}{293757}
\definecolor{footergray}{gray}{0.65}

\geometry{a4paper,left=3cm,right=3cm,top=3cm,bottom=3cm}
\renewcommand{\familydefault}{\sfdefault}
\setcounter{secnumdepth}{3}

\titleformat{\section}
  {\color{nato}\Huge\bfseries}{\thesection.}{1em}{}
\titlespacing*{\section}{0pt}{30pt}{12pt}

\titleformat{\subsection}
  {\color{nato}\LARGE\bfseries}{\thesubsection.}{1em}{}
\titlespacing*{\subsection}{0pt}{20pt}{8pt}

\titleformat{\subsubsection}
  {\color{nato}\Large\bfseries}{\thesubsubsection.}{1em}{}
\titlespacing*{\subsubsection}{0pt}{10pt}{6pt}

\fancyhf{}
\fancyhead[R]{\small\color{footergray} 29 червня 2026}
\fancyfoot[L]{\small\color{footergray} Техноробочий проект ERP/1. Модуль IAS}
\fancyfoot[R]{\small\color{footergray}\thepage}

\newcommand\HUGE[1]{\fontsize{40}{40}\selectfont #1}

\renewcommand{\headrulewidth}{0pt}
\renewcommand{\footrulewidth}{0.4pt}
\renewcommand{\footrule}{\color{footergray}\hrule width \textwidth height 0.4pt}

\begin{document}

\pagestyle{fancy}

\begin{titlepage}
  \AddToShipoutPictureBG*{\AtPageLowerLeft{\color{nato}\rule{\paperwidth}{\paperheight}}}
  \color{white}\thispagestyle{empty}\vspace*{3cm}\centering
  {\Huge  \textbf{ERP/1: Авторизація} \par} \vspace{1.5cm}
  {\HUGE  \textbf{Техноробочий проект} \par} \vspace{1.5cm}
  {\LARGE \textbf{на створення та впровадження модуля «Служба ідентифікації та авторизації (IAS)» як засобу інформатизації} \par} \vspace{1.5cm}
  {\large Згідно з ДСТУ 34.201-89 та Постановою КМУ № 205 від 21.02.2025 \par}
  {\large Формалізація вимог за стандартом ISO/IEC/IEEE 29148:2018 \par}
  \vfill
  {\large 2026 \copyright\ Системи Електронного Урядування \par}
\end{titlepage}

\newpage

\begin{center}
  {\Huge\color{nato}\bfseries Зміст\par}
\end{center}
\vspace{40pt}
\tableofcontents

\newpage

\section{Загальні відомості про засіб інформатизації}

\subsection{Найменування та підстави для розробки}
\textbf{Найменування засобу інформатизації:} Модуль «Служба ідентифікації та авторизації (IAS)» (далі --- Модуль IAS / Система IAS) у складі інформаційної системи управління підприємством ERP/1.

\textbf{Підстави для розробки:}
\begin{enumerate}
    \item Закон України «Про захист інформації в інформаційно-телекомунікаційних системах».
    \item Постанова Кабінету Міністрів України від 29.03.2006 № 373 «Про затвердження Правил забезпечення захисту інформації в інформаційних, телекомунікаційних та інформаційно-телекомунікаційних системах».
    \item Постанова Кабінету Міністрів України від 21.02.2025 № 205 «Про затвердження Порядку використання засобів інформатизації».
    \item Угода про асоціацію між Україною та Європейським Союзом (в частині транскордонної взаємодії електронних довірчих послуг).
    \item Технічне завдання на побудову Zencrypted інфраструктури безпеки.
\end{enumerate}

\subsection{Статуси опису реалізації}
Щоб не змішувати фактичну реалізацію, частково завершені можливості, цільову архітектуру та нормативні наміри, у цьому документі застосовуються такі статуси:
\begin{longtable}{p{0.24\textwidth}p{0.68\textwidth}}
\toprule
\textbf{Статус} & \textbf{Зміст} \\
\midrule
Реалізовано & Підтверджено поточним кодом, стійкими структурами даних і автоматизованими тестами IAS. \\
Частково реалізовано & Реалізовано лише для визначеного сценарію або з тимчасовими технічними залежностями. \\
Цільовий стан & Запланована архітектура або функція наступних ітерацій; не трактується як наявна у поточному середовищі виконання. \\
Нормативна вимога & Вимога законодавства, стандарту, сертифікації або організаційної політики, що потребує окремих доказів реалізації та приймання. \\
\bottomrule
\end{longtable}

\subsection{Простежуваність вимог та відповідність нормативній базі}
Техноробочий проєкт деталізує технічні вимоги з \texttt{ias.tex} і має зберігати зв'язок між вимогою, фактичним модулем, стійкою структурою даних, автоматизованим тестом та статусом реалізації. Посилання на неіснуючі модулі, демонстраційний код або майбутню архітектуру не повинні подаватися як опис поточного стану.

На цьому етапі зафіксовано такі межі простежуваності:
\begin{itemize}
    \item \textbf{Керування доменними об'єктами:} користувачі, пристрої, сертифікати, профілі безпеки, політики та VPN-сервіси зберігаються у KVS/Mnesia через доменні модулі IAS. Статус: Реалізовано.
    \item \textbf{Майстер налаштування VPN:} поточний життєвий цикл виконується спеціалізованими модулями IAS з перевірками переходів, готовності матеріалів, доставки та відновлення. Статус: Реалізовано.
    \item \textbf{Випуск сертифікатів через CA:} поточний CMP-шлях використовує зовнішній командний інтерфейс OpenSSL і є перехідним. Повне усунення запуску OpenSSL із середовища виконання IAS є цільовим станом.
    \item \textbf{Керування процесами через BPE/BPMN:} майбутнє перенесення керування станами й переходами на \texttt{synrc/bpmn}/BPE є цільовим станом; поточний майстер поки не виконується рушієм BPE.
    \item \textbf{EST, CMC, CRL/OCSP, повний ABAC, КСЗІ та національні криптографічні алгоритми:} позначаються лише як цільові або нормативні вимоги до появи окремих реалізаційних і приймальних доказів.
\end{itemize}

\subsection{Призначення та цілі створення засобу}
Модуль IAS є центральною адміністративно-координаційною службою для керування цифровими ідентичностями, пристроями, сертифікатами, профілями безпеки, політиками та VPN-сервісами. IAS формує бажаний стан і передає керовані конфігурації до VPN, але не бере участі у пересиланні VPN-трафіку.

\textbf{Поточні та цільові цілі:}
\begin{itemize}
    \item \textbf{Реалізовано:} стійкий реєстр доменних об'єктів у KVS/Mnesia, граф зв'язків, перевірка CSR і X.509, покрокове налаштування пристрою, передавання конфігурацій до VPN через Erlang RPC, ревізії, контрольні суми, журналювання доставки, звіряння стану та відновлення після збоїв.
    \item \textbf{Частково реалізовано:} випуск сертифікатів через CA за CMP із зовнішнім командним інтерфейсом OpenSSL як перехідною залежністю.
    \item \textbf{Цільовий стан:} керований клієнт або адаптер CA без запуску зовнішнього OpenSSL у IAS, а також перенесення керування життєвим циклом майстра на \texttt{synrc/bpmn}/BPE зі збереженням наявних доменних модулів і перевіреної поведінки.
    \item \textbf{Нормативна вимога:} окреме підтвердження вимог КСЗІ, КЕП, криптографії за ДСТУ, захисту персональних даних та інших нормативних профілів.
\end{itemize}

\subsection{Послідовність виконання наступних робіт}
Календарні строки не фіксуються без окремої оцінки обсягу, ресурсів і критеріїв приймання. Роботи виконуються у такій залежній послідовності:
\begin{enumerate}
    \item \textbf{P0 --- синхронізація документації:} привести \texttt{ias.tex} і \texttt{ias\_pro.tex} у відповідність до поточного коду та однакових статусів реалізації.
    \item \textbf{P1 --- усунення OpenSSL із IAS:} прибрати запуск зовнішнього командного інтерфейсу OpenSSL із середовища виконання IAS, зберігши CMP або інший погоджений контракт через керований клієнт чи адаптер CA.
    \item \textbf{P2 --- стабілізація наскрізного сценарію:} прогнати й стабілізувати повний життєвий цикл майстра налаштування VPN від користувача й пристрою до передавання конфігурації, повторних спроб і відновлення після збоїв.
    \item \textbf{P3 --- опис поточного процесу:} зафіксувати фактичну FSM/BPMN-модель станів, переходів, умов, помилок і завершення без розбіжностей із кодом.
    \item \textbf{P4 --- перенесення керування на BPE:} поетапно зробити BPE джерелом стану й переходів, залишивши майстер користувацьким інтерфейсом над процесом та зберігши доменні сервісні модулі.
    \item \textbf{P5 --- повторна синхронізація:} узгодити технічні вимоги, техноробочий проєкт, випробування та профілі безпеки після реалізації цільових змін.
\end{enumerate}

\newpage
\section{Відомості про робочий процес та умови експлуатації}

\subsection{Опис бізнес-процесів (BPMN / FSM)}
Поточні процеси IAS реалізовано спеціалізованими модулями та стійкими записами стану, а не рушієм BPE. Наведені нижче процесні описи мають бути перевірені за фактичним кодом у наступних ітераціях цього документа: поточний стан слід описувати як AS IS FSM, а BPE/BPMN --- як цільову модель перенесення керування станами й переходами. Фрагменти BPE DSL у поточній редакції не є доказом наявної реалізації.

\subsubsection{Процес реєстрації пристрою та імпорту OVPN (bpe\_ovpn\_import.erl)}
Цей процес відповідає за початкове імпортування раніше створених OpenVPN конфігурацій та створення відповідних реєстрових сутностей в IAS.
\begin{enumerate}
    \item \textbf{StartImport:} Початок процесу оператором безпеки.
    \item \textbf{ParseConfig:} Сервісна задача, яка парсить .ovpn та створює структуру вилучених параметрів.
    \item \textbf{VerifyCertificates:} Перевірка валідності імпортованого сертифіката (підпис, термін дії, ланцюжок довіри).
    \item \textbf{ApproveImportPlan:} Користувацька задача затвердження згенерованого плану імпорту оператором.
    \item \textbf{CreateGraphObjects:} Сервісна задача, яка записує об'єкти Device, Certificate, VPN Service та зв'язки між ними в KVS.
    \item \textbf{ImportCompleted:} Завершення процесу.
\end{enumerate}

\subsubsection{Процес випуску та сертифікації CMP (bpe\_cmp\_enrollment.erl)}
Процес забезпечує автоматизоване отримання сертифіката від CA за допомогою CSR, згенерованого безпосередньо клієнтським пристроєм.
\begin{enumerate}
    \item \textbf{StartEnrollment:} Ініціація пристроєм запиту на випуск.
    \item \textbf{GenerateKeyCSRPlan:} Створення в IAS плану генерації ключів та CSR на пристрої.
    \item \textbf{AwaitCSRUpload:} Очікування завантаження файлу CSR з пристрою.
    \item \textbf{ValidateCSR:} Перевірка підпису та параметрів CSR.
    \item \textbf{SubmitToCMP:} Відправка CSR до CA через протокол CMP.
    \item \textbf{VerifyCertificate:} Отримання випущеного сертифіката, перевірка публічного ключа проти CSR, запис у сховище та оновлення статусу пристрою.
    \item \textbf{EnrollmentCompleted:} Завершення випуску.
\end{enumerate}

\subsubsection{Процес авторизації та оцінки довіри (bpe\_trust\_evaluation.erl)}
Процес динамічної оцінки права пристрою чи користувача на підключення.
\begin{enumerate}
    \item \textbf{StartEvaluation:} Виклик процесу під час handshake пристрою або планової перевірки.
    \item \textbf{CheckExpiry:} Перевірка дати закінчення дії сертифіката.
    \item \textbf{CheckCRL:} Запит та перевірка серійного номера за списками відкликання (CRL/OCSP).
    \item \textbf{EvaluateABACRules:} Оцінка контекстних правил (роль, ЄДРПОУ, час, робоче місце).
    \item \textbf{UpdateTrustStatus:} Запис оновленого статусу довіри пристрою (Ready або Blocked).
    \item \textbf{EvaluationCompleted:} Завершення оцінки.
\end{enumerate}

\subsubsection{Процес провіженінгу VPN (bpe\_vpn\_provisioning.erl)}
Процес транзакційної доставки налаштувань на шлюзи VPN та контролю узгодження.
\begin{enumerate}
    \item \textbf{StartProvisioning:} Ініціація оновлення конфігурації.
    \item \textbf{BuildCommand:} Генерація канонічної команди провіженінгу з інкрементним номером ревізії.
    \item \textbf{DeliverToVPN:} Виклик віддаленої процедури RPC на шлюзі VPN.
    \item \textbf{AwaitCompletionNotice:} Очікування повідомлення від vpn\_event\_bus про успішне застосування змін на шлюзі.
    \item \textbf{VerifyDataplane:} Перевірка проходження тестового пакету через зашифрований тунель.
    \item \textbf{ProvisioningCompleted:} Завершення провіженінгу.
\end{enumerate}

\subsection{Опис ролей та прав доступу (ABAC)}
Доступ до адміністративних функцій IAS обмежується на базі ABAC. Приклад правила авторизації:
\begin{verbatim}
allow(Operator, Action, Object) ->
    OperatorRole = maps:get(role, Operator),
    OperatorEDRPOU = maps:get(edrpou, Operator),
    ObjectSensitivity = maps:get(sensitivity_level, Object),
    
    IsSecurityAdmin = (OperatorRole == security_admin),
    MatchesEDRPOU = (OperatorEDRPOU == maps:get(edrpou, Object, undefined)),
    
    case {IsSecurityAdmin, ObjectSensitivity} of
        {true, _} -> true;
        {false, confidential} -> MatchesEDRPOU;
        {false, public} -> true;
        _ -> false
    end.
\end{verbatim}

\subsection{Умови експлуатації та системне середовище}
Система IAS повинна розгортатися у кластерному середовищі:
\begin{itemize}
    \item \textbf{Операційна система:} Linux (Ubuntu Server LTS, Rocky Linux) або macOS (для розробки).
    \item \textbf{Платформа:} Erlang/OTP версії 28+.
    \item \textbf{Служби взаємодії:} Розподілені Erlang-ноди, з'єднані за допомогою epmd та TLS mesh.
    \item \textbf{Вимоги до СУБД:} Mnesia, інтегрована через інтерфейс KVS.
\end{itemize}

\newpage
\section{Інформаційне та програмне забезпечення}

\subsection{Інформаційне забезпечення (Схеми та Дані)}

\subsubsection{Визначення структур даних (Erlang Records)}
Всі ключові сутності описуються у вигляді Erlang рекорді:

\begin{verbatim}
-record(ias_domain_object, {
    key,                       % {Kind, ObjectId}
    schema_version = 1,        % Версія схеми даних
    kind,                      % user | device | certificate | policy
    object_id,                 % Унікальний UUID
    payload = #{},             % Мапа з атрибутами об'єкта
    revision = 1,              % Номер поточної ревізії
    created_at = 0,            % Timestamp створення
    updated_at = 0             % Timestamp оновлення
}).

-record(ias_certificate_material_record, {
    key,                       % {SubjectKind, SubjectId}
    schema_version = 1,
    subject_kind,              % user | device | ca
    subject_id,                % UUID суб'єкта
    material_type,             % public_cert | cert_chain | signature
    encoding = pem,            % pem | der
    source,                    % cmp | est | manual
    fingerprint_sha256,        % SHA256 відбиток сертифіката
    body_envelope = #{},       % Мапа, що містить PEM-тіло
    revision = 1,
    created_at = 0,
    updated_at = 0,
    expires_at = undefined     % Термін дії сертифіката
}).

-record(ias_vpn_device_state, {
    device_id,                 % UUID пристрою (PK)
    schema_version = 2,
    revision = 0,              % Актуальна ревізія на шлюзі
    command_digest = undefined,% Хеш останньої відправленої команди
    canonical_command = #{},   % Тіло сформованої команди
    binding = #{},             % Параметри мережевої прив'язки
    lifecycle_state = unbound, % unbound | pending | active | degraded
    last_decommission = undefined,
    decommission_history = [],
    decommissioned_at = undefined,
    updated_at = 0
}).

-record(ias_csr_enrollment_record, {
    csr_fingerprint,           % Хеш CSR (PK)
    schema_version = 1,
    status = submitted,        % submitted | processing | signed | rejected
    retryable = false,         % Чи можна повторити спробу
    payload = #{},             % Мапа з параметрами запиту та CSR PEM
    revision = 1,
    created_at = 0,
    updated_at = 0
}).

-record(ias_vpn_reconciliation_incident, {
    device_id,                 % UUID пристрою
    schema_version = 1,
    kind,                      % orphan | stale_revision | fingerprint_mismatch
    reason,                    % Текстовий опис розбіжності
    token,                     % Унікальний токен інциденту
    status = open,             % open | acknowledged | resolved
    snapshot = #{},            % Порівняльний знімок станів IAS та VPN
    first_seen = 0,
    last_seen = 0,
    occurrences = 1,
    acknowledged_by = undefined,
    acknowledged_note = undefined,
    acknowledged_at = undefined,
    resolved_by = undefined,
    resolved_note = undefined,
    resolved_at = undefined,
    updated_at = 0
}).
\end{verbatim}

\subsubsection{Опис довідників та реєстрових сутностей}
Опис полів реєстрових сутностей (таких як \texttt{User}, \texttt{Device}, \texttt{Certificate}) відповідає таксономії, зазначеній у Технічних вимогах. Всі записи зберігаються у таблиці Mnesia \texttt{kvs}, яка містить записи типу \texttt{ias\_domain\_object} та інші допоміжні рекорди.

\subsubsection{Алгоритми оцінки статусу довіри (Trust Evaluation)}
Приклад визначення статусу пристрою на основі його криптографічного стану:
\begin{verbatim}
evaluate_device_trust(DeviceState, CertRecord, CurrentTime) ->
    case CertRecord#ias_certificate_material_record.expires_at of
        Expires when CurrentTime > Expires ->
            {blocked, expired};
        _ ->
            case check_crl_status(CertRecord#ias_certificate_material_record.fingerprint_sha256) of
                revoked -> {blocked, revoked};
                ok ->
                    case DeviceState#ias_vpn_device_state.lifecycle_state of
                        degraded -> {degraded, connection_timeout};
                        active -> {ready, ok};
                        unbound -> {incomplete, no_binding}
                    end
            end
    end.
\end{verbatim}

\subsection{Програмне забезпечення (Реалізація)}

\subsubsection{bpe\_ovpn\_import.erl (DSL процесу імпорту)}
\begin{verbatim}
-module(bpe_ovpn_import).
-include("bpe.hrl").
-compile(export_all).

def() ->
    #process{
        name = 'OVPN Profile Import',
        beginEvent = 'StartImport',
        endEvent = 'ImportCompleted',
        tasks = [
            #beginEvent{name='StartImport'},
            #serviceTask{name='ParseConfig', module=?MODULE},
            #serviceTask{name='VerifyCertificates', module=?MODULE},
            #userTask{name='ApproveImportPlan', module=?MODULE},
            #serviceTask{name='CreateGraphObjects', module=?MODULE},
            #endEvent{name='ImportCompleted'}
        ],
        flows = [
            #sequenceFlow{name='1', source='StartImport', target='ParseConfig'},
            #sequenceFlow{name='2', source='ParseConfig', target='VerifyCertificates'},
            #sequenceFlow{name='3', source='VerifyCertificates', 
                          target='ApproveImportPlan', condition="certs_ok"},
            #sequenceFlow{name='4', source='VerifyCertificates', 
                          target='StartImport', condition="certs_invalid"},
            #sequenceFlow{name='5', source='ApproveImportPlan', target='CreateGraphObjects'},
            #sequenceFlow{name='6', source='CreateGraphObjects', target='ImportCompleted'}
        ],
        roles = [security_admin, operator]
    }.

action({serviceTask, 'ParseConfig'}, Proc) ->
    % Логіка парсингу OVPN файлу
    {reply, Proc, "certs_ok"};
action({serviceTask, 'CreateGraphObjects'}, Proc) ->
    % Створення сутностей в KVS
    {reply, Proc};
action(_, Proc) -> {reply, Proc}.
\end{verbatim}

\subsubsection{bpe\_cmp\_enrollment.erl (DSL процесу CMP випуску)}
\begin{verbatim}
-module(bpe_cmp_enrollment).
-include("bpe.hrl").
-compile(export_all).

def() ->
    #process{
        name = 'CA CMP Enrollment',
        beginEvent = 'StartEnrollment',
        endEvent = 'EnrollmentCompleted',
        tasks = [
            #beginEvent{name='StartEnrollment'},
            #serviceTask{name='GenerateKeyCSRPlan', module=?MODULE},
            #userTask{name='AwaitCSRUpload', module=?MODULE},
            #serviceTask{name='ValidateCSR', module=?MODULE},
            #serviceTask{name='SubmitToCMP', module=?MODULE},
            #serviceTask{name='VerifyCertificate', module=?MODULE},
            #endEvent{name='EnrollmentCompleted'}
        ],
        flows = [
            #sequenceFlow{name='1', source='StartEnrollment', target='GenerateKeyCSRPlan'},
            #sequenceFlow{name='2', source='GenerateKeyCSRPlan', target='AwaitCSRUpload'},
            #sequenceFlow{name='3', source='AwaitCSRUpload', target='ValidateCSR'},
            #sequenceFlow{name='4', source='ValidateCSR', target='SubmitToCMP', condition="csr_valid"},
            #sequenceFlow{name='5', source='ValidateCSR', target='AwaitCSRUpload', condition="csr_invalid"},
            #sequenceFlow{name='6', source='SubmitToCMP', target='VerifyCertificate', condition="enrollment_ok"},
            #sequenceFlow{name='7', source='SubmitToCMP', target='StartEnrollment', condition="enrollment_failed"},
            #sequenceFlow{name='8', source='VerifyCertificate', target='EnrollmentCompleted'}
        ],
        roles = [operator, device]
    }.

action({serviceTask, 'ValidateCSR'}, Proc) ->
    {reply, Proc, "csr_valid"};
action({serviceTask, 'SubmitToCMP'}, Proc) ->
    {reply, Proc, "enrollment_ok"};
action(_, Proc) -> {reply, Proc}.
\end{verbatim}

\subsubsection{bpe\_trust\_evaluation.erl (DSL процесу перевірки довіри)}
\begin{verbatim}
-module(bpe_trust_evaluation).
-include("bpe.hrl").
-compile(export_all).

def() ->
    #process{
        name = 'Device Trust Evaluation',
        beginEvent = 'StartEvaluation',
        endEvent = 'EvaluationCompleted',
        tasks = [
            #beginEvent{name='StartEvaluation'},
            #serviceTask{name='CheckExpiry', module=?MODULE},
            #serviceTask{name='CheckCRL', module=?MODULE},
            #serviceTask{name='EvaluateABACRules', module=?MODULE},
            #serviceTask{name='UpdateTrustStatus', module=?MODULE},
            #endEvent{name='EvaluationCompleted'}
        ],
        flows = [
            #sequenceFlow{name='1', source='StartEvaluation', target='CheckExpiry'},
            #sequenceFlow{name='2', source='CheckExpiry', target='CheckCRL', condition="not_expired"},
            #sequenceFlow{name='3', source='CheckExpiry', target='UpdateTrustStatus', condition="expired"},
            #sequenceFlow{name='4', source='CheckCRL', target='EvaluateABACRules', condition="not_revoked"},
            #sequenceFlow{name='5', source='CheckCRL', target='UpdateTrustStatus', condition="revoked"},
            #sequenceFlow{name='6', source='EvaluateABACRules', target='UpdateTrustStatus'},
            #sequenceFlow{name='7', source='UpdateTrustStatus', target='EvaluationCompleted'}
        ],
        roles = [security_admin, system]
    }.

action({serviceTask, 'CheckExpiry'}, Proc) ->
    {reply, Proc, "not_expired"};
action({serviceTask, 'CheckCRL'}, Proc) ->
    {reply, Proc, "not_revoked"};
action(_, Proc) -> {reply, Proc}.
\end{verbatim}

\subsubsection{bpe\_vpn\_provisioning.erl (DSL процесу провіженінгу VPN)}
\begin{verbatim}
-module(bpe_vpn_provisioning).
-include("bpe.hrl").
-compile(export_all).

def() ->
    #process{
        name = 'VPN Command Provisioning',
        beginEvent = 'StartProvisioning',
        endEvent = 'ProvisioningCompleted',
        tasks = [
            #beginEvent{name='StartProvisioning'},
            #serviceTask{name='BuildCommand', module=?MODULE},
            #serviceTask{name='DeliverToVPN', module=?MODULE},
            #userTask{name='AwaitCompletionNotice', module=?MODULE},
            #serviceTask{name='VerifyDataplane', module=?MODULE},
            #endEvent{name='ProvisioningCompleted'}
        ],
        flows = [
            #sequenceFlow{name='1', source='StartProvisioning', target='BuildCommand'},
            #sequenceFlow{name='2', source='BuildCommand', target='DeliverToVPN'},
            #sequenceFlow{name='3', source='DeliverToVPN', target='AwaitCompletionNotice', condition="delivery_ok"},
            #sequenceFlow{name='4', source='DeliverToVPN', target='StartProvisioning', condition="delivery_failed"},
            #sequenceFlow{name='5', source='AwaitCompletionNotice', target='VerifyDataplane'},
            #sequenceFlow{name='6', source='VerifyDataplane', target='ProvisioningCompleted'}
        ],
        roles = [operator, system]
    }.

action({serviceTask, 'BuildCommand'}, Proc) ->
    {reply, Proc};
action({serviceTask, 'DeliverToVPN'}, Proc) ->
    {reply, Proc, "delivery_ok"};
action(_, Proc) -> {reply, Proc}.
\end{verbatim}

\subsubsection{ias\_validator.erl (Модуль валідації)}
\begin{verbatim}
-module(ias_validator).
-export([validate_csr/1, validate_chain/2, verify_fingerprint/2]).

validate_csr(CSR_PEM) ->
    % Виклик OpenSSL або вбудованого ASN.1 розшифровувача для перевірки підпису CSR
    case public_key:pem_decode(CSR_PEM) of
        [_] -> ok;
        _   -> {error, invalid_csr_format}
    end.

validate_chain(Cert_PEM, CA_Chain_PEM) ->
    % Криптографічна перевірка ланцюжка сертифікатів CA
    ok.

verify_fingerprint(Cert_PEM, SHA256) ->
    % Порівняння відбитку SHA-256
    Calculated = crypto:hash(sha256, Cert_PEM),
    case Calculated of
        SHA256 -> ok;
        _      -> {error, fingerprint_mismatch}
    end.
\end{verbatim}

\newpage
\section{Вимоги до засобу за ISO/IEC/IEEE 29148}

\subsection{Функціональні вимоги}
Система повинна забезпечувати повнофункціональне керування ідентичностями, валідацію сертифікатів, оновлення конфігурацій шлюзів VPN без втрати з'єднань користувачів та виявлення невідповідностей стану мережі.

\subsection{Вимоги до інтерфейсів та дизайну}
\begin{itemize}
    \item Веб-інтерфейс адміністративної консолі повинен базуватися на реактивному фреймворку N2O/NITRO.
    \item Відповідність вимогам WCAG 2.1 рівню AA для доступності інтерфейсів.
    \item Візуальний дизайн повинен наслідувати фірмовий стиль Zencrypted з підтримкою Sleek Dark Mode.
\end{itemize}

\subsection{Вимоги до безпеки та захисту інформації}
\begin{itemize}
    \item Шифрування каналів зв'язку за допомогою TLS v1.3.
    \item Підтримка КЕП CAdES-X Long та алгоритму ДСТУ 4145-2002.
    \item Ізоляція ключів: приватний ключ ніколи не записується і не передається в IAS, залишаючись виключно на кінцевому пристрої.
    \item Атестація комплексу за вимогами ДССЗІ України на рівень захисту Г3.
\end{itemize}

\subsection{Вимоги до продуктивності та надійності}
\begin{itemize}
    \item Час аптайму системи не менше 99.99\% (High Availability).
    \item Час відновлення (RTO) після повного перезапуску ноди не більше 5 секунд завдяки KVS/Mnesia rehydration.
    \item Здатність обробляти до 1000 запитів авторизації на секунду на одну ноду.
\end{itemize}

\subsection{Вимоги до логування та аудиту}
\begin{itemize}
    \item Всі операції провіженінгу, випуску та зміни статусів довіри повинні підписуватися цифровим підписом ноди та записуватися в стійкий лог аудиту.
    \item Заборона модифікації або видалення записів журналу аудиту (immutable audit log).
\end{itemize}

\subsection{Адміністративні та юридичні вимоги}
\begin{itemize}
    \item Всі компоненти поставляються з ліцензіями відкритого вихідного коду (MIT/Apache 2.0).
    \item Модуль IAS повинен інтегруватися з наявною інфраструктурою підприємства без потреби закупівлі додаткового платного ПЗ.
\end{itemize}

\end{document}
